\chapter{Controllo remoto, comunicazione App e miglioramenti di stabilita}

\section{Obiettivo del capitolo}
Questo capitolo documenta tutte le aggiunte introdotte per la comunicazione remota con app mobile: architettura software, protocollo comandi, strategia di test, ottimizzazioni per i limiti hardware e problemi riscontrati in Wokwi con relative soluzioni.

\section{Architettura della comunicazione remota}
La comunicazione e stata progettata a strati:

\begin{itemize}
  \item \textbf{Astrazione trasporto}: \texttt{bluetooth\_connection.h} definisce l'interfaccia comune.
  \item \textbf{Trasporto BLE concreto}: \texttt{ble\_link.h} implementa l'interfaccia per Bluetooth Low Energy.
  \item \textbf{Logica protocollo}: \texttt{command\_processor.h} interpreta i comandi, valida i payload e aggiorna lo stato.
  \item \textbf{Orchestrazione}: \texttt{remote\_gateway.h} gestisce seriale, polling BLE e pubblicazione periodica dello stato.
\end{itemize}

Questa struttura semplifica manutenzione e sostituzione del componente di comunicazione.

\section{Protocollo utilizzato dall'app}
Il protocollo e testuale e pensato per integrazione semplice con App Inventor.

\subsection{Comandi base}
\begin{itemize}
  \item \texttt{PING} $\rightarrow$ \texttt{PONG}
  \item \texttt{VER} $\rightarrow$ \texttt{VER:1}
  \item \texttt{GET:STATE}
  \item \texttt{GET:MODE}
  \item \texttt{MANUAL:ON} / \texttt{MANUAL:OFF}
\end{itemize}

\subsection{Comandi setpoint}
\begin{itemize}
  \item \texttt{SET:TEMP:<valore>} con range $[15.0,30.0]$
  \item \texttt{SET:HUM:<valore>} con range $[40.0,80.0]$
  \item \texttt{SET:LUX:<valore>} con range $[50,1200]$
  \item \texttt{SET:PEOPLE:<valore>} con range $[0,\mathrm{maxPeople}]$
\end{itemize}

\subsection{Comandi override manuale}
In modalita \texttt{MANUAL} l'app puo bypassare la logica automatica:
\begin{itemize}
  \item Servo: \texttt{SERVO:OPEN}, \texttt{SERVO:CLOSE}, \texttt{SERVO:ANGLE:<0..180>}
  \item Attuatori: \texttt{ACT:GREEN:*}, \texttt{ACT:RED:*}, \texttt{ACT:HEAT:*}, \texttt{ACT:COOL:*}, \texttt{ACT:HUMIDIFIER:*}
  \item PWM luci: \texttt{ACT:PLAFONIERE:PWM:<0..255>}
\end{itemize}

\subsection{Formato payload stato}
Esempio:
\begin{lstlisting}[style=cppstyle]
STATE:T=21.8;H=58.0;L=170;P=3;S=90;M=MANUAL;TT=22.5;TH=60.0;TL=180
\end{lstlisting}

\subsection{Gestione errori}
Per robustezza lato app sono previsti errori espliciti:
\begin{itemize}
  \item \texttt{ERR:EMPTY}, \texttt{ERR:UNKNOWN}, \texttt{ERR:MODE}, \texttt{ERR:SERVO}
  \item Errori di range: \texttt{ERR:RANGE:*}
  \item Errori di formato: \texttt{ERR:FORMAT:*}
  \item Protezione overflow seriale: \texttt{ERR:TOO\_LONG}
\end{itemize}

\section{Strategia di validazione e test}
La suite test e stata estesa con casi dedicati al controllo remoto:

\begin{itemize}
  \item correttezza protocollo (comandi validi e invalidi),
  \item valori limite per tutti i range,
  \item payload numerici malformati,
  \item transizioni AUTO/MANUAL,
  \item input seriale frammentato,
  \item input seriale troppo lungo,
  \item trasporto disconnesso,
  \item coda comandi su cicli multipli.
\end{itemize}

Nei test viene usato un trasporto finto per validare la logica anche senza hardware BLE reale.

\section{Ottimizzazioni software per stabilita}
Sono state introdotte diverse ottimizzazioni per ridurre rischi di instabilita:

\begin{itemize}
  \item riduzione footprint RAM con tipi interi piu piccoli per pin e contatori,
  \item rimozione di comportamento bloccante nel tornello e chiusura servo non bloccante,
  \item riduzione churn dinamico delle stringhe nella generazione stato,
  \item parsing numerico rigoroso per evitare conversioni non sicure,
  \item separazione output build per target diversi,
  \item workflow dedicati per simulazione e hardware reale.
\end{itemize}

Queste modifiche migliorano prevedibilita runtime e riducono problemi in simulazione/deploy.

\section{Problemi riscontrati in Wokwi e soluzioni}

\subsection{Problema A: mismatch board/firmware}
Se la board simulata e il firmware compilato non coincidono, la simulazione risulta non funzionante.

\textbf{Soluzione}: output separati:
\begin{itemize}
  \item Wokwi in \texttt{build-wokwi/}
  \item hardware UNO R4 in \texttt{build-r4/}
\end{itemize}

\subsection{Problema B: chiave duplicata in \texttt{wokwi.toml}}
Una chiave \texttt{elf} duplicata causava errore di parsing TOML.

\textbf{Soluzione}: mantenere una sola entry \texttt{elf} e validare il file.

\subsection{Problema C: servo non visibile in movimento}
Dopo la rimozione dei delay bloccanti, il servo poteva aprire/chiudere troppo rapidamente nello stesso ciclo.

\textbf{Soluzione}: introduzione chiusura temporizzata non bloccante e trigger su fronte pulsante.

\section{Workflow build completo (comando unico)}
Il progetto supporta un comando unico che esegue test completi e doppia compilazione:

\begin{lstlisting}[style=cppstyle]
make full-test
\end{lstlisting}

Questo comando esegue:
\begin{enumerate}
  \item test completi,
  \item compilazione firmware Wokwi,
  \item compilazione firmware UNO R4 WiFi.
\end{enumerate}

\section{Test app senza hardware}
E stato aggiunto un mock server locale per iniziare test App Inventor senza scheda fisica:

\begin{itemize}
  \item script: \texttt{tools/mock\_arduino\_server.py}
  \item comando: \texttt{make mock-server}
\end{itemize}

Con questo approccio puoi validare UI, parsing e flussi comandi prima dei test BLE reali.

\section{Nota finale di deploy}
Prima del rilascio finale e comunque necessario validare BLE su hardware UNO R4 WiFi reale, per verificare radio, latenza e comportamento in condizioni reali.
